% Book pp. 97-102 (PDF pp. 98-103)
\subsection{The Rational Zero Theorem}

The Rational Zero Theorem helps us identify potential rational zeros of a
polynomial by looking at the factors of two specific parts of the polynomial:
\begin{enumerate}
  \item Factors of the Constant Term: This is the term without a variable (the last
  number in the polynomial when arranged in descending powers).
  \item Factors of the Leading Coefficient: This is the coefficient (the number in
  front of the term) of the term with the highest power of $x$.
\end{enumerate}

\noindent\textbf{\textit{Steps to Find Possible Rational Zeros}}
\begin{enumerate}
  \item Identify the Constant Term and Leading Coefficient.

  For example, in the polynomial $4x^3 + 7x^2 + 2x + 5$:

  The constant term is 5.

  The leading coefficient is 4.
  \item List the Factors.

  Factors of the Constant Term (5): These are: $\pm 1$, $\pm 5$

  Factors of the Leading Coefficient (4): These are $\pm 1$, $\pm 2$, $\pm 4$.
  \item Create Possible Rational Zeros.

  To find the possible rational zeros, take each factor of the constant term and
  divide it by each factor of the leading coefficient.

  Possible rational zeros for our example would be:

  From $\pm 1$: $\frac{1}{1}, -\frac{1}{1}, \frac{1}{2}, -\frac{1}{2}, \frac{1}{4},
  \frac{-1}{4}, = \pm 1, \pm\frac{1}{2}, \pm\frac{1}{4}$,

  From $\pm 5$: $\frac{5}{1} - \frac{5}{1}, \frac{5}{2}, -\frac{5}{2}, \frac{5}{4},
  \frac{-5}{4}, = \pm 5, \pm\frac{5}{2}, \pm\frac{5}{4}$
  \item List of Possible Rational Zeros:

  Combining all these, the possible rational zeros for our polynomial are:

  $\pm 1, \pm 5, \pm\frac{1}{2}, \pm\frac{1}{4}, , \pm\frac{5}{2}, \pm\frac{5}{4}$
\end{enumerate}
Now that we have a list of possible rational zeros, we can use the Factor Theorem
to test each one.
\begin{enumerate}
  \item[a.] For each rational number c from the list, substitute it into the polynomial
  f(x)
  \item[b.] If f($c$) $=$ 0, then c is a zero of the polynomial.
\end{enumerate}

\begin{example}{Example 3.4}
In groups or individually, use the rational zeros theorem to list all the possible
rational zeros.
\begin{enumerate}
  \item[a.] $6x^3 + 15x^2 - 8x + 2$
  \item[b.] $x^3 + x^2 - 2x - 5$
  \item[c.] $3x^3 + x^2 + x - 5$
  \item[d.] $6x^3 - 4x^2 + 17$
\end{enumerate}
\booknote[S3-02]{the solution below works parts b and c for different
polynomials, $x^3 + x^2 - 2x - 8$ and $9x^3 + x^2 + x - 10$, from those printed
in the question.}

\solution
\begin{enumerate}
  \item[a.] $6x^3 + 15x^2 - 8x + 2$

  The constant term is 2.

  The leading coefficient is 6.

  Factors of the Constant Term (2): These are.$\pm 1$, $\pm 2$

  Factors of the Leading Coefficient (6): These are $\pm 1$, $\pm 2$, $\pm 3$, $\pm 6$.

  Possible Rational Zeros:

  From $\pm 1$: $\frac{1}{1}, -\frac{1}{1}, \frac{1}{2}, -\frac{1}{2}, \frac{1}{3},
  \frac{-1}{3}, \frac{1}{6}, -\frac{1}{6} = \pm 1, \pm\frac{1}{2}, \pm\frac{1}{3},
  \pm\frac{1}{6}$,

  From $\pm 2$: $\frac{2}{2}, -\frac{2}{2}, \frac{2}{3}, \frac{-2}{3}, \frac{2}{6},
  -\frac{2}{6} = \pm 1, \pm\frac{2}{3}, \pm\frac{1}{6}$,
  \booknote[S3-03]{the book's ``From $\pm 2$'' line omits $\frac{2}{1}$ and ends
  in $\pm\frac{1}{6}$; dividing $\pm 2$ by $1, 2, 3, 6$ gives $\pm 2, \pm 1,
  \pm\frac{2}{3}, \pm\frac{1}{3}$, and no combined list is given for this part.}

  \item[b.] $x^3 + x^2 - 2x - 8$

  The constant term is $-8$.

  The leading coefficient is 1.

  Factors of the Constant Term (-8): These are.$\pm 1$, $\pm 2$, $\pm 4$, $\pm 8$

  Factors of the Leading Coefficient (1): These are $\pm 1$.

  Possible Rational Zeros:

  From $\pm 1$: $\frac{1}{1}, -\frac{1}{1} = \pm 1$

  From $\pm 2$: , $\frac{2}{1}, -\frac{2}{1} = \pm 2$

  From $\pm 4$: $\frac{4}{1}, -\frac{4}{1}, = \pm 4$

  From $\pm 8$: , $\frac{8}{1}, -\frac{8}{1} = \pm 8$,,

  Possible Rational Zeros (combined): $\pm 1, \pm 2, \pm 4, \pm 8$

  \item[c.] $9x^3 + x^2 + x - 10$

  The constant term is $-10$.

  The leading coefficient is 9.

  Factors of the Constant Term (-10): These are.$\pm 1$, $\pm 2$, $\pm 5$, $\pm 10$

  Factors of the Leading Coefficient (9): These are $\pm 1$, $\pm 3$, $\pm 9$.

  Possible Rational Zeros:

  From $\pm 1$: $\frac{1}{1}, -\frac{1}{1}, \frac{1}{3}, -\frac{1}{3}, \frac{1}{9},
  -\frac{1}{9} = \pm 1, \pm\frac{1}{3}, \pm\frac{1}{9}$

  From $\pm 2$: $\frac{2}{1}, -\frac{2}{1}, \frac{2}{3}, -\frac{2}{3}, \frac{2}{9},
  -\frac{2}{9} = \pm 2, \pm\frac{2}{3}, \pm\frac{2}{9}$

  From $\pm 4$: $\frac{5}{1}, -\frac{5}{1}, \frac{5}{3}, -\frac{5}{3}, \frac{5}{9},
  -\frac{5}{9} = \pm 5, \pm\frac{5}{3}, \pm\frac{5}{9}$
  \booknote[S3-04]{the book heads this line ``From $\pm 4$''; it should be
  ``From $\pm 5$'', as the fractions show.}

  From $\pm 10$: , $\frac{10}{1}, -\frac{10}{1}, \frac{10}{3}, -\frac{10}{3},
  \frac{10}{9}, -\frac{10}{9} = \pm 10, \pm\frac{10}{3}, \pm\frac{10}{9}$

  Possible Rational Zeros (combined): $\pm 1, \pm 2, \pm 5, \pm 10\ \pm\frac{1}{3},
  \pm\frac{1}{9}, \pm\frac{2}{3}, \pm\frac{2}{9}, \pm\frac{5}{3}, \pm\frac{5}{9},
  \pm\frac{10}{3}, \pm\frac{10}{9}$

  \item[d.] $6x^3 - 4x^2 + 17$

  The constant term is 17.

  The leading coefficient is 6.

  Factors of the Constant Term (17): These are.$\pm 1$, $\pm 17$

  Factors of the Leading Coefficient (6): These are $\pm 1$, $\pm 2$, $\pm 3$, $\pm 6$.

  Possible Rational Zeros\textbf{:}

  From $\pm 1$: $\frac{1}{1}, -\frac{1}{1}, \frac{1}{2}, -\frac{1}{2}, \frac{1}{3},
  -\frac{1}{3}, \frac{1}{6}, -\frac{1}{6} = \pm 1, \pm\frac{1}{2}, \pm\frac{1}{3},
  \pm\frac{1}{6}$

  From $\pm 17$: $\frac{17}{1}, -\frac{17}{1}, \frac{17}{2}, -\frac{17}{2},
  \frac{17}{3}, -\frac{17}{3}, \frac{17}{6}, -\frac{17}{6} = = \pm 17,
  \pm\frac{17}{2}, \pm\frac{17}{3}, \pm\frac{17}{6}$

  Possible Rational Zeros (combined): $\pm 1, \pm 17, \pm\frac{1}{2}, \pm\frac{1}{3},
  \pm\frac{1}{6}, \pm\frac{17}{2}, \pm\frac{17}{3}, \pm\frac{17}{6}$
\end{enumerate}
\end{example}

\subsection{Using the rational zero theorem to find the zeros of a polynomial function}

\noindent\textit{Steps to Follow:}
\begin{enumerate}
  \item Identify Possible Rational Zeros

  Use the Rational Zero theorem to list all possible rational zeros. These are
  of the form $p$ and $q$, where $p$ is a factor of the constant term (the last term of
  the polynomial) and $q$ is a factor of the leading coefficient (the coefficient of
  the term with the highest degree).
  \item Substitute and Evaluate
  \begin{enumerate}
    \item[a.] Substitute each possible rational zero into the polynomial function.
    \item[b.] Calculate the value of the polynomial for each substitution.
  \end{enumerate}
  \item Determine Actual Zeros
  \begin{enumerate}
    \item[a.] Identify which substitutions result in a value of zero.
    \item[b.] The values for which the polynomial equals zero are the rational zeros
    of the function.
  \end{enumerate}
\end{enumerate}

\begin{example}{Example 3.5}
Given the polynomial $3x^3 + 8x^2 + 7x + 2$
\begin{enumerate}
  \item List the factors of the constant term and leading coefficient.
  \item Determine possible rational zeros
  \item Substitute each value into the polynomial to find which ones yield zero.
\end{enumerate}

\solution
The constant term is 2.

The leading coefficient is 3.

Factors of the Constant Term (2): These are.$\pm 1$, $\pm 2$

Factors of the Leading Coefficient (3): These are $\pm 1$, $\pm 3$.

Possible Rational Zeros:

From $\pm 1$: $\frac{1}{1}, -\frac{1}{1}, \frac{1}{3}, -\frac{1}{3} = \pm 1,
\pm\frac{1}{3}$

From $\pm 2$: $\frac{2}{1}, -\frac{2}{1}, \frac{2}{3}, -\frac{2}{3}, = = \pm 2,
\pm\frac{2}{3}$

Possible Rational Zeros (combined): $\pm 1, \pm 2, \pm\frac{1}{3}, \pm\frac{2}{3}$

Substitute and evaluate:

$3x^3 + 8x^2 + 7x + 2$

For $x = \pm 1$:

$x = 1$

$3(1)^3 + 8(1)^2 + 7(1) + 2$

$3 + 8 + 7 + 2 = 20$

$\boldsymbol{x} = -1$

$3(-1)^3 + 8(-1)^2 + 7(-1) + 2$

$-3 + 8 - 7 + 2 = 0$

For $x = \pm 2$

$x = 2$

$3(2)^3 + 8(2)^2 + 7(2) + 2$

$24 + 32 + 14 + 2 = 71$

$x = -2$

$3(-2)^3 + 8(-2)^2 + 7(-2) + 2$

$-24 + 32 - 14 + 2 = -4$

For $\pm\frac{1}{3}$

$x = \frac{1}{3}$

$3\left(\frac{1}{3}\right)^3 + 8\left(\frac{1}{3}\right)^2 + 7\left(\frac{1}{3}\right) + 2$

$\frac{1}{9} + \frac{8}{9} + \frac{7}{3} + 2 = \frac{16}{3}$

$x = -\frac{1}{3}$

$3\left(-\frac{1}{3}\right)^3 + 8\left(-\frac{1}{3}\right)^2 + 7\left(-\frac{1}{3}\right) + 2$

$-\frac{1}{9} + \frac{8}{9} - \frac{7}{3} + 2 = \frac{4}{9}$

For $\pm\frac{2}{3}$

$x = \frac{2}{3}$

$3\left(\frac{2}{3}\right)^3 + 8\left(\frac{2}{3}\right)^2 + 7\left(\frac{2}{3}\right) + 2$

$\frac{8}{9} + \frac{32}{9} + \frac{14}{3} + 2 = \frac{100}{9}$

$x = -\frac{2}{3}$

$3\left(-\frac{2}{3}\right)^3 + 8\left(-\frac{2}{3}\right)^2 + 7\left(-\frac{2}{3}\right) + 2$

$-\frac{8}{9} + \frac{32}{9} - \frac{14}{3} + 2 = 0$

The values $1, 2, -2, \frac{1}{3}, -\frac{1}{3}, \frac{2}{3}$ did not yield zero and are
not zeros of the function.

However, $-1$ and $-\frac{2}{3}$, yielded zero and therefore represent the zeros of
the function
\end{example}
